% Capítulo textual do trabalho.

\chapter{Introdução}
\label{cap:introducao}

A introdução é uma seção fundamental em qualquer trabalho acadêmico, pois estabelece as bases para a pesquisa a ser realizada. Nela, o autor deve apresentar o tema que será investigado, situando-o no contexto atual e destacando sua relevância para o campo de estudo. Conforme destacado por \citeonline{severino}, a introdução deve fornecer uma visão geral clara e concisa, permitindo ao leitor compreender as motivações e os objetivos da pesquisa.

A primeira etapa na elaboração da introdução é a contextualização do tema. É necessário que o autor apresente uma descrição detalhada do cenário em que o problema de pesquisa está inserido, evidenciando a importância do assunto em questão. De acordo com \citeonline{lakatos}, essa contextualização é essencial para que o leitor perceba a pertinência da investigação, além de indicar as lacunas existentes na literatura que o trabalho busca preencher.

Em seguida, a introdução deve incluir a formulação do problema de pesquisa. Esta etapa é crucial, uma vez que a clareza e especificidade da pergunta norteadora influenciam diretamente o desenvolvimento do trabalho. Como mencionado por \citeonline{creswell}, a formulação adequada do problema estabelece os limites da pesquisa e direciona as investigações subsequentes. Portanto, é imperativo que o problema seja apresentado de maneira explícita e compreensível.

Outro aspecto relevante da introdução é a definição dos objetivos da pesquisa. Os objetivos devem ser claramente enunciados, subdividindo-se em um objetivo geral e objetivos específicos. O objetivo geral reflete a meta maior do estudo, enquanto os objetivos específicos detalham as etapas necessárias para alcançar esse propósito. Essa estrutura ajuda a guiar o leitor ao longo da pesquisa, fornecendo um panorama sobre o que se pretende investigar~\cite{gil}.

Adicionalmente, a introdução deve abordar de forma breve a metodologia que será utilizada ao longo do trabalho. Embora a descrição completa da metodologia seja mais adequada em uma seção específica, oferecer uma visão geral na introdução ajuda o leitor a compreender como o autor planeja abordar o problema de pesquisa. Segundo \citeonline{koche}, a metodologia é a estratégia que o pesquisador utiliza para alcançar os objetivos e responder ao problema formulado.

A introdução de um trabalho acadêmico deve ser redigida com clareza e objetividade, abordando os elementos centrais da pesquisa, como a contextualização do tema, a formulação do problema, a justificativa, os objetivos e a metodologia. Além disso, a introdução pode fazer uso de elementos não textuais, como ilustrações e tabelas, para facilitar a compreensão da metodologia que será desenvolvida. No entanto, a introdução deve ser linear e não subdividida, ou seja, a divisão em subseções sugere detalhamento, o que deve ser reservado para as seções específicas.

Finalmente, a introdução deve concluir com uma breve descrição da estrutura do trabalho, indicando como os capítulos estão organizados. Essa abordagem permite ao leitor ter uma compreensão clara do fluxo do texto e dos tópicos que serão discutidos ao longo da pesquisa. A seguir, apresenta-se a estrutura do trabalho.

Este trabalho está organizado como segue.
O Capítulo~\ref{cap:organizacao} descreve a estrutura típica de trabalhos acadêmicos de graduação.
O Capítulo~\ref{cap:citacao} aborda as normas para formatação e uso de citações, distinguindo entre citação direta, indireta e de fontes.
O Capítulo~\ref{cap:ilustracoes} orienta sobre a inserção de ilustrações e gráficos, enquanto o Capítulo~\ref{cap:tabelas} esclarece a diferença entre tabelas e quadros, com exemplos práticos.
O Capítulo~\ref{cap:codigos} trata da apresentação de códigos e algoritmos.
O Capítulo~\ref{cap:equacoes} explora o uso de equações, e o Capítulo~\ref{cap:alineas} discute as regras para o uso de alíneas.
Por fim, o Capítulo~\ref{cap:cronograma} trata da formatação de cronogramas, e o Capítulo~\ref{cap:conclusao} apresenta as considerações finais.
